\newcommand{\docshorttitle}{NTC LODGING SUPPL.\ AUTH.\ --- \S 425.17(c) CROSS-HEARING (CGC-25-631801)}
\newcommand{\doctitleblock}{NOTICE OF LODGING SUPPLEMENTAL AUTHORITY (CODE CIV.\ PROC., \S\S~425.16, 425.17)}
\newcommand{\docsubblock}{(Cal.\ Rules of Court, rule 3.1306(c); lodge only; not an amendment to any pleading)}
% Pleading-paper preamble for APRIL-24-2026 cover sheets and oral-argument
% outlines (CGC-25-631801 and CGC-25-631802). Matches house style from
% APRIL-23-SAC-DEFENSE/court-pdfs/REPLY-SAC-PREAMBLE.tex, simplified for
% single-to-three-page artifacts that do not require TOC/TOA infrastructure.
%
% Cross-hearing caption (CAPTION-801/802.tex): use \CaptionDocTitleBlock,
% \CaptionDocSubBlock, and the crosshearingsschedule environment (defined below)
% so long \doctitleblock text and hearing dates stay inside the 3in minipage.
\documentclass{article}
\usepackage[utf8]{inputenc}
\usepackage{graphicx}
\usepackage{geometry}
    \geometry{top=1in, bottom=2.2in, left=1.0in, right=1.0in, textheight=336pt}
    \setlength{\footskip}{70pt}
\usepackage{ulem}
\usepackage{fancyhdr}
    \renewcommand{\headrulewidth}{0pt}
    \renewcommand{\footrulewidth}{0.4pt}
\usepackage{parskip}
    \setlength{\parskip}{0mm}
\usepackage{quoting}
    \quotingsetup{leftmargin=0.5in, rightmargin=0.5in, vskip=-1.5mm}
    \makeatletter
        \g@addto@macro\quoting\singlespacing
        \g@addto@macro\quoting{\vspace{-2mm}}
    \makeatother
\usepackage{mathptmx}
\usepackage{amssymb}
\renewcommand{\rmdefault}{ptm}
\renewcommand{\normalsize}{\fontsize{12}{12}\selectfont}
\newcommand*{\ptm}{\fontfamily{ptm} \selectfont \fontsize{12}{12}\selectfont}
\usepackage{setspace}
    \doublespacing
\newcommand{\tab}{\hspace*{.0in}}
\newenvironment{tightcenter}{%
    \setlength\topsep{0pt}
    \setlength\parskip{0pt}
    \begin{center}
}{%
    \end{center}
}
\usepackage{tikz}
\usetikzlibrary{calc}
% Pleading-paper vertical double rule at the left margin.
\AddToHook{shipout/background}{%
    \begin{tikzpicture}[remember picture, overlay]
        \draw[line width=0.5pt] ($(current page.north west) + (0.75in,0)$) -- ($(current page.south west) + (0.75in,0)$);
        \draw[line width=0.5pt] ($(current page.north west) + (0.80in,0)$) -- ($(current page.south west) + (0.80in,0)$);
    \end{tikzpicture}%
}
\usepackage[left, pagewise]{lineno}
\setlength{\linenumbersep}{0.35in}
\renewcommand{\linenumberfont}{\normalfont\small}
\usepackage{titlesec}
\usepackage[hidelinks]{hyperref}
\usepackage{bookmark}
\usepackage{textcomp}
\usepackage{longtable}
\usepackage{booktabs}
\usepackage{array}
\usepackage{multirow}
\usepackage{calc}
% --- Cross-hearing caption cover (3in minipage): compact title + 3-column hearing table ---
% Use on first page only; keeps long \doctitleblock / \docsubblock on one sheet with party caption.
\newcommand{\CaptionDocTitleBlock}{%
  \begingroup
  \fontsize{10}{11.5}\selectfont
  \bfseries
  \raggedright
  \setlength{\parskip}{0pt}%
  \noindent\doctitleblock\par
  \endgroup
}
\newcommand{\CaptionDocSubBlock}{%
  \begingroup
  \footnotesize
  \raggedright
  \setlength{\parskip}{0pt}%
  \noindent\docsubblock\par
  \endgroup
}
% #1 = department number only (e.g.\ 301 or 302) for the line ``Dept.\ #1''.
\newenvironment{crosshearingsschedule}[1]{%
  \par\begingroup
  \footnotesize
  \setlength{\tabcolsep}{2.2pt}%
  \renewcommand{\arraystretch}{1.04}%
  \setlength{\parskip}{0pt}%
  \setlength{\parindent}{0pt}%
  \noindent\textbf{Cross-noticed for pending defense hearings (this case, Dept.\ #1):}\par\vspace{2pt}%
  \begin{tabular}{@{}>{\raggedright\arraybackslash}p{0.11\linewidth}>{\raggedright\arraybackslash}p{0.34\linewidth}>{\raggedright\arraybackslash}p{0.53\linewidth}@{}}%
  \textbf{\#} & \textbf{Motion} & \textbf{Date / dept / judge} \\ \hline
}{%
  \end{tabular}\par\endgroup\smallskip
}
% Pandoc pipe-table column widths use \real{#}.
\newcommand{\real}[1]{\pgfmathparse{#1}\pgfmathresult}
\titleformat{\section}{\fontsize{13}{13}\selectfont\normalfont\bfseries\uppercase}{}{0em}{}
\titleformat{\subsection}{\fontsize{12}{12}\selectfont\normalfont\bfseries}{}{0em}{}
\titleformat{\subsubsection}{\normalfont\bfseries}{}{0em}{}
\titlespacing*{\section}{0pt}{6pt}{2pt}
\titlespacing*{\subsection}{0pt}{4pt}{2pt}
\titlespacing*{\subsubsection}{0pt}{2pt}{2pt}
\newcommand{\calrules}[1]{Cal. Rules of Court, rule #1}
\newcommand{\ccp}[1]{Code Civ. Proc., \textsection~#1}
\newcommand{\evidcode}[1]{Evid. Code, \textsection~#1}
\providecommand{\tightlist}{%
  \setlength{\itemsep}{0pt}\setlength{\parskip}{0pt}}
% Footer: page N of total + document short-title.
\usepackage{lastpage}
\pagestyle{fancy}
\fancyhf{}
\fancyfoot[C]{\thepage{} of \pageref{LastPage} \\[1pt] \footnotesize \docshorttitle}
\fancyfoot[R]{}

\begin{document}
% Cross-hearing caption for CGC-25-631801 (Dept. 301, Hon. Van Aken).
\nolinenumbers
\begin{singlespace}
\noindent Franciscus Dylan Rosario \\
7624 Melody Drive \\
Rohnert Park, CA 94928 \\
Tel: 650.488.7510 \\
E: soltrinox@gmail.com \\
Plaintiff, In Pro Per
\end{singlespace}
\vspace*{9mm}
\begin{tightcenter}
\textbf{SUPERIOR COURT OF THE STATE OF CALIFORNIA} \\
\textbf{COUNTY OF SAN FRANCISCO}
\end{tightcenter}
\vspace*{3mm}
\begin{minipage}[t]{3in}
    \raggedright
    \textbf{FRANCISCUS DYLAN ROSARIO,} \\ \textbf{PLAINTIFF}, \\
    \hspace{1cm} v. \\
    \small
    Subhi Abdelhalim; CSAA Insurance Exchange; Phillips, Spallas \& Angstadt LLP; Priya D. Navaratnasingham; Alberto Reyna; Michael R. Chambers; Carbone, Smith \& Koyama; and Does 1 through 50, inclusive;\\
    \normalsize
    \textbf{DEFENDANTS}
    \makebox[3in]{\hrulefill}
\end{minipage}%
\hspace{3mm}%
\begin{minipage}[t]{0.5pt}
    \vspace{0pt}
    \rule{0.5pt}{2.42in}
\end{minipage}%
\hspace{3mm}%
\begin{minipage}[t]{3in}
    \raggedright
    \noindent\textbf{Case No.: CGC-25-631801}\par\smallskip
    \CaptionDocTitleBlock\smallskip
    \CaptionDocSubBlock\smallskip
    \begin{singlespace}
    \setlength{\parindent}{0pt}
    \begin{crosshearingsschedule}{301}
    (1) & MFL / leave (SAC) & \textbf{Apr.\ 30, 2026} \\
    (2) & Anti-SLAPP (\ccp{425.16}) & \textbf{May 12, 2026, 9:30 a.m.};\ Dept.\ 301;\ Hon.\ Christine Van Aken \\
    (3) & Strike (\ccp{436}) & \textbf{May 22, 2026};\ Dept.\ 301;\ Hon.\ Christine Van Aken \\
    (4) & Demurrer & Hrng.\ date \textbf{[TBD]} (reserved) \\
    \end{crosshearingsschedule}
    \noindent Complaint filed: Dec.\ 23, 2025 \\
    \noindent FAC filed: Feb.\ 25, 2026 \\
    \end{singlespace}
\end{minipage}
\vspace*{8mm}
\newpage
\linenumbers


\section*{NOTICE OF LODGING}

TO ALL PARTIES AND THEIR ATTORNEYS OF RECORD:

PLEASE TAKE NOTICE that Plaintiff Franciscus Dylan Rosario hereby \textbf{lodges supplemental authority} for the Court's consideration at \textbf{each} of the pending defense-motion hearings cross-listed on the caption page: MFL/SAC (Apr.\ 30, 2026); special motion to strike (May 12, 2026); \ccp{436} motion to strike (May 22, 2026); and demurrer (date reserved). The authorities below bear on whether \ccp{425.17(c)} and related threshold doctrines are \textbf{common spine} issues across those motions.

\subsection*{I. Authorities Lodged}

\begin{enumerate}
\tightlist
\item \textit{Simpson Strong-Tie Co.\ v.\ Gore} (2010) 49 Cal.4th 12 (framework for \ccp{425.17(c)}; categorical bar when elements satisfied).
\item \textit{JAMS, Inc.\ v.\ Superior Court} (2016) 1 Cal.App.5th 984 (\ccp{425.17(c)} applied to commercial actor making representations about its own services).
\item \textit{Stewart v.\ Rolling Stone LLC} (2010) 181 Cal.App.4th 664 (\ccp{425.17(c)} applied to statements integrated with commercial output).
\item \textit{Flatley v.\ Mauro} (2006) 39 Cal.4th 299 (illegality as a matter of law removes assertedly protected conduct from \ccp{425.16} gateway).
\item \textit{Lefebvre v.\ Lefebvre} (2011) 199 Cal.App.4th 696 (reaffirming \textit{Flatley} on uncontroverted record).
\item \textit{Club Members for an Honest Election v.\ Sierra Club} (2008) 45 Cal.4th 309 (threshold construction of \ccp{425.17} exemptions before \ccp{425.16}).
\item \textit{Action Apartment Ass'n v.\ City of Santa Monica} (2007) 41 Cal.4th 1232 (\ccp{47}(b) coextensive with \ccp{425.16}(e) for certain deceit predicates).
\item \textit{Rusheen v.\ Cohen} (2006) 37 Cal.4th 1048 (non-communicative conduct outside \ccp{425.16}(e)).
\item \textit{City of Cotati v.\ Cashman} (2002) 29 Cal.4th 69 (demurrer pleading standards; rhetorical merits labels).
\item \textit{Park v.\ Board of Trustees} (2017) 2 Cal.5th 1057; \textit{Baral v.\ Schnitt} (2016) 1 Cal.5th 376 (anti-SLAPP step separation).
\item \textit{Egan v.\ Mutual of Omaha Ins.\ Co.} (1979) 24 Cal.3d 809; \textit{Wilson v.\ 21st Century Ins.\ Co.} (2007) 42 Cal.4th 713 (insurer investigation duties).
\item \textit{Zhang v.\ Superior Court} (2013) 57 Cal.4th 364 (UCL predicate and Insurance Code standards).
\item Plaintiff's concurrently lodged Omnibus Supplemental Memorandum of Black-Letter Contradictions (systematizing defense positions against cited authorities).
\end{enumerate}

\subsection*{II. Common Defense Spine Across All Four Motions}

The operative complaint and incorporated papers explain why \ccp{425.17(c)} and privilege sequencing are not one-off labels tied only to the anti-SLAPP calendar. The same commercial-insurer and panel-counsel posture appears in demurrer Ground B framing, in the \ccp{436} motion papers that invoke ``absolute'' litigation-privilege themes, and in the MFL/SAC sequencing where defendants invoke \ccp{425.16} policy rhetoric. Plaintiff does not amend any pleading here; this lodging supplies supplemental authority and does not enlarge the merits briefing record. Physical copies of all lodged materials accompany this Notice pursuant to \calrules{3.1306(c)}.

\subsection*{III. Why \ccp{425.17(c)} Reaches All Four Motions}

\textbf{Threshold first.} Under \textit{Club Members}, \ccp{425.17} issues are structurally prior where raised. \textbf{\ccp{47}(b) coextensive strand.} Under \textit{Action Apartment}, certain deceit channels align with the \ccp{425.16}(e) inquiry. \textbf{\ccp{436} frame.} A \ccp{436} motion to strike is not a free-standing substitute for an anti-SLAPP special motion; where defendants import anti-SLAPP-style ``everything is petitioning'' arguments into a \ccp{436} attack, \textit{PH II, Inc.\ v.\ Superior Court} (1995) 33 Cal.App.4th 1680, 1682--83 (backdoor demurrer misuse) supplies a strike-specific guardrail Plaintiff lodges for the May 22, 2026 hearing. \textbf{\textit{Simmons} class cases.} Amendment futility rhetoric tied to \ccp{425.16} (often associated with \textit{Simmons v.\ Allstate Ins.\ Co.} (2001) 92 Cal.App.4th 1068) does not erase threshold \ccp{425.17} sequencing where the Court is asked to rule on leave and related papers on the Apr.\ 30, 2026 calendar; Plaintiff addresses that interaction through \textit{JKC3H8 v.\ Colton} (2013) 221 Cal.App.4th 468 and \textit{Nguyen-Lam v.\ Cao} (2009) 171 Cal.App.4th 871.

\subsection*{IV. \ccp{425.17}(e) --- No Interlocutory Stay or Appeal}

\ccp{425.17}(e) provides that an order granting or denying a motion under \ccp{425.17} is not subject to an interlocutory appeal, stay, or writ to the extent it is based on a \ccp{425.17} exemption. Plaintiff incorporates by reference his lodged Interlocutory Appeal Bar brief on \ccp{904.1(a)(13)} and related writ constraints.

\subsection*{V. \ccp{425.17}(d) --- No Carve-Out Available on This Record}

\ccp{425.17}(d) limits certain ``public interest'' carve-out arguments for particular categories of defendants. On the facts Plaintiff has pleaded about commercial insurance claim handling and paid defense counsel, defendants cannot reposition this dispute as a ``media / internet / large nonprofit'' \ccp{425.17}(d) case to sidestep \ccp{425.17}(c). The lodged authorities in Section I therefore remain relevant at each hearing where defendants invoke broad \ccp{425.16} immunity tropes.

\subsection*{VI. Procedural Posture}

This lodging does not amend or expand the briefing record. Physical copies of all lodged materials are provided herewith pursuant to \calrules{3.1306(c)}. An accompanying RJN seeks judicial notice of the operative statute and regulation texts, the full text of which is reproduced therein.

\subsection*{VII. Relief Requested}

Plaintiff requests that the Court (a) accept this Notice for the register, (b) consider the lodged authorities at \textbf{each} hearing listed on the caption page, and (c) take judicial notice as requested in the accompanying RJN.

\input{SIGBLOCK.tex}
\end{document}
